\scanpage{38}
domaine analytique complexe). Ici
$X=\mathbb C$ est topologiquement trivial, ce qui
peut s'exprimer ainsi par le fait que
son type d'homotopie est complètement
décrit par celui du topos $\unsure{\mathcal F|\mathfrak F}$ des
systèmes (17). En effet, on aura
\[
(22)\qquad \unsure{\mathcal F|\mathfrak F}\simeq\widehat C
\]
où $C$ (dans le cas général \deleted{du}\inserted{où dans le} diagramme
(19) \deleted{\unsure{avec $\pi_1(Y)$ réduit à $\{1\}$}}
\inserted{où $\pi_1(Y)=\{1\}$},
\inserted{$D\Rightarrow\pi_1(\mathcal V^*_{Y,X})\xrightarrow{\sim}\pi_1(X^{\unsure{*|\ell}})$ iso)}
\deleted{\unsure{alors}}) est \unsure{défini|définie} comme la catégorie dont
l'ens des objets est l'ens \deleted{\unsure{somme}}\inserted{$\{a,b\}$,}
\deleted{\unsure{réunion} \illegible{3} \unsure{des trois}}
\deleted{\unsure{flèches de sous-catégorie}}
\deleted{\unsure{des ens d'objets} des groupoïdes, composants
\inserted{\unsure{figurer en la} \deleted{\unsure{une}} catégorie}
de (8), et dont les flèches sont, en plus
des flèches de ces trois groupoïdes, les
flèches $\deleted{\unsure{p}}$ d'un objet $a$ de $\inserted{A=}\Pi_1\mathcal V^*_{Y,X}$
vers un objet soit de $B=\Pi_1\deleted{\unsure{\mathcal V_{Y,X}|\mathcal V^*_{Y,X}}}\inserted{\unsure{Y}}$, soit de
$C=\Pi_1\deleted{\unsure{\mathcal V_{Y,X}|\mathcal V^*_{Y,X}}}\inserted{X^*}$, \unsure{définies|définis} comme étant
les flèches :
\[
\deleted{\unsure{\varphi}}\quad p_*(a)\longrightarrow b\quad\text{dans }B
\quad(\text{resp. }j_*(a)\longrightarrow c\quad\text{dans }C).
\]
Dans le cas \deleted{\unsure{actuel}}, présent, les groupoïdes
$A,B,C$ et les morphismes entre eux}
% Pale lower-left foliation, in another hand; not the author's page number.
\margin{37}
